\begin{table}[!htbp]
\centering
\footnotesize
\renewcommand{\arraystretch}{1.15}
\caption{\textbf{Similarity of Strategic-Game Contexts to Receiver-Type Vignettes}}
\label{tab:similarity_receivers}
\begin{tabular}{l cc}
\toprule
Context & \multicolumn{2}{c}{Graded similarity (0--100)} \\
\midrule
\multicolumn{3}{l}{\emph{Ultimatum game: Accepting vs.\ Rejecting}} \\
Control & 47.0 & 58.8 \\
Market & 50.0 & 56.7 \\
Market $-$ Control & +3.0 & -2.2 \\
\midrule
\multicolumn{3}{l}{\emph{Trust game: Returning vs.\ Keeping}} \\
Control & 62.5 & 42.5 \\
Market & 61.0 & 43.5 \\
Market $-$ Control & -1.5 & +1.0 \\
\bottomrule
\end{tabular}
\begin{flushleft}
\footnotesize Notes: Mean 0--100 ratings of how similar the counterpart evoked by each context is to the receiver-type vignettes (two vignettes per type, three permuted-label conversations; six ratings per cell), from the round-2 headline rater under the frozen protocol of Appendix~\ref{app:similarity_vignettes}. Contexts are the verbatim Control and Market instructions of the ultimatum and trust games. The forced 100-point splits over the same vignettes are in Table~\ref{tab:receiver_splits}.
\end{flushleft}
\end{table}
